%
% chapters/chapter02/main.tex
% Chapter 2: Model development and experimental setup
%

\chapter{Model development and experimental setup}
\label{chap:setup}

\section{Data collection and preparation}
\label{sec:data}

\subsection{Description of the data collection process (transport complaint records from City Transportation Systems, \textasciitilde3\,000 records)}
\label{sec:data_collection}

The dataset was provided directly by City Transportation Systems (CTS), the public transport operator of Astana, Kazakhstan. The raw corpus consists of approximately 3\,000 complaint records collected through operational channels including a mobile application, a web portal, and call-centre transcriptions. Each record contains a free-text description field (\foreignlanguage{russian}{\textit{описание}}) in which passengers describe transportation issues: missed stops, route deviations, driver misconduct, and payment failures.

Records were collected over a twelve-month operational period. The distribution of complaint types reflects the uneven nature of public transport incidents: vehicle identification complaints (route number, bus ID) are most common, while geographic intersection complaints (where exactly a problem occurred) are the rarest but most actionable for dispatch. No personally identifiable information was included; all records were anonymised by CTS prior to transfer.

\subsection{Description of the annotation process (manual BIO labelling in Label Studio, six entity types)}
\label{sec:annotation}

Manual annotation was performed using \textit{Label Studio}, an open-source data-labelling platform. Three annotators with native-level proficiency in both Kazakh and Russian independently labelled each complaint record with the six entity types defined in Section~\ref{sec:geo_entities}. Annotation guidelines were developed iteratively, with inter-annotator disagreements resolved through adjudication to produce a gold-standard label set.

Tokens were labelled according to the BIO scheme. Annotators applied span-level labels directly on raw tokenised text, selecting the exact substring corresponding to each entity. Adjudication followed a majority-vote protocol: any span where at least two of three annotators agreed was accepted as gold standard; spans with three-way disagreement were reviewed by the project supervisor.

The final annotated dataset comprises 2\,750 complaint records exported in JSONL format with BIO token-level labels.

\subsection{Description of preprocessing methods (tokenization, subword alignment, handling Russian--Kazakh code-switching)}
\label{sec:preprocessing}

Raw complaint text was preprocessed through the following pipeline:

\begin{enumerate}
  \item \textbf{Unicode normalisation.} Cyrillic characters were normalised to eliminate look-alike homoglyphs (Latin-Cyrillic confusables introduced by keyboard layouts).
  \item \textbf{Whitespace tokenisation.} Records were split on whitespace and punctuation to produce word-level tokens. Punctuation tokens were retained as separate items.
  \item \textbf{Subword alignment.} For transformer models, word-level BIO labels were aligned to subword tokens produced by the model-specific tokeniser (SentencePiece BPE for XLM-R and mE5; WordPiece for ruBERT). Only the first subword of each word received the original label; subsequent subword pieces were masked with $-100$ (excluded from loss computation).
  \item \textbf{Code-switching handling.} No explicit language identification or separation was performed. Multilingual models (XLM-R, mE5) process mixed-language tokens natively within their shared vocabulary. The LLM pipeline (Qwen2.5) handled code-switching at the prompt level through bilingual few-shot examples.
\end{enumerate}

\subsection{Dividing the data into training and test samples (70/30 stratified split, 1\,925 / 825 records)}
\label{sec:split}

The annotated dataset was split into training and test sets using stratified sampling to maintain entity type proportions across splits:

\begin{itemize}
  \item \textbf{Training set:} 1\,925 samples (70\,\%), used for model fine-tuning.
  \item \textbf{Test set:} 825 samples (30\,\%), used exclusively for evaluation.
\end{itemize}

Stratified sampling was applied at the record level, using the presence or absence of each entity type as a stratification key. Table~\ref{tab:dataset} presents the resulting dataset statistics by entity type.

\begin{table}[H]
\centering
\caption{Dataset statistics by entity type.}
\label{tab:dataset}
\renewcommand{\arraystretch}{1.3}
\begin{tabular}{lrrrr}
\toprule
\textbf{Entity Type} & \textbf{Train Tokens} & \textbf{Test Tokens} & \textbf{Total} & \textbf{Avg./Record} \\
\midrule
\texttt{ROUTE\_NUM}          & 2{,}641 & 1{,}131 & 3{,}772 & 1.37 \\
\texttt{BUS\_ID}             &   798 &   342 & 1{,}140 & 0.41 \\
\texttt{PLATE\_NUM}          &   623 &   267 &   890 & 0.32 \\
\texttt{STOP\_NAME}          & 3{,}901 & 1{,}672 & 5{,}573 & 2.03 \\
\texttt{STREET\_NAME}        &   287 &   123 &   410 & 0.15 \\
\texttt{STREET\_INTERSECTION}&    52 &    22 &    74 & 0.03 \\
\midrule
\textbf{Total}               & 8{,}302 & 3{,}557 & 11{,}859 & -- \\
\bottomrule
\end{tabular}
\end{table}

\section{Description of the machine learning models used}
\label{sec:methods}

Four machine learning approaches were implemented. The BIO label set comprises 13 entity labels ($B$- and $I$- for each of the 6 entity types) plus the $O$ class, yielding $K = 14$ output classes. All fine-tuning experiments were conducted on an NVIDIA Tesla T4 GPU (16\,GB VRAM).

\subsection{Model 1: XLM-RoBERTa (multilingual transformer)}
\label{sec:xlmr}

XLM-RoBERTa (\texttt{xlm-roberta-base})~\cite{conneau2020xlmr} is a multilingual transformer encoder pre-trained on 2.5~TB of CommonCrawl data across 100 languages using a masked language modelling objective. Its cross-lingual transfer capabilities and coverage of both Kazakh and Russian make it well-suited for the code-switching characteristics of the CTS corpus.

\textbf{Architecture.} The base encoder (12 transformer layers, 768 hidden dimensions, 12 attention heads, $\approx$125M parameters) is extended with a token classification head, a linear layer $W \in \mathbb{R}^{768 \times K}$ projecting each contextual embedding to class logits. Token probabilities are computed via the softmax function:
\begin{equation}
  p(y_i = c \mid \mathbf{x}_i) = \frac{\exp(z_{i,c})}{\sum_{k=1}^{K} \exp(z_{i,k})}
  \label{eq:softmax}
\end{equation}
where $\mathbf{x}_i$ is the contextual embedding of token $i$, $z_{i,c}$ is the raw logit for class $c$, and $K=14$.

\textbf{Loss function.} Training uses weighted cross-entropy to address class imbalance:
\begin{equation}
  \mathcal{L} = -\frac{1}{N} \sum_{i=1}^{N} w_{y_i} \log p(y_i \mid \mathbf{x}_i), \qquad w_c = \frac{N_{\text{total}}}{K \cdot N_c}
  \label{eq:wce}
\end{equation}
where $N$ is the number of valid (non-padding) tokens, $N_{\text{total}}$ is the total number of labelled tokens, and $N_c$ is the frequency of class $c$. The \texttt{STREET\_NAME} class receives an additional manual upweight factor of $1.6$.

\textbf{Learning rate schedule.} A linear decay is applied:
\begin{equation}
  \text{lr}(t) = \text{lr}_{\max} \cdot \left(1 - \frac{t}{T}\right)
\end{equation}
where $\text{lr}_{\max} = 2 \times 10^{-5}$, $t$ is the current training step, and $T$ is the total number of steps.

\textbf{Training configuration:} 10 epochs, batch size 8, AdamW optimiser (weight decay 0.01), SentencePiece BPE tokeniser (vocabulary size 250\,000).

\subsection{Model 2: ruBERT (DeepPavlov, Russian BERT)}
\label{sec:rubert}

ruBERT (\texttt{DeepPavlov/rubert-base-cased})~\cite{kuratov2019rubert} is a BERT-base model pre-trained on Russian Wikipedia and news corpora ($\approx$180M parameters). It serves as a monolingual Russian-language baseline, providing a direct measure of the performance penalty incurred by ignoring the Kazakh portion of the code-switching corpus.

\textbf{Token label alignment.} Sub-word tokenisation is handled via offset mapping:
\begin{equation}
  L'_i = \begin{cases} L_i & \text{if } \text{offset}_i \neq (0,0) \\ -100 & \text{otherwise} \end{cases}
\end{equation}
where $L_i$ is the original BIO label and $-100$ masks special/padding tokens from the loss.

\textbf{Prediction.} Token label assignment follows:
\begin{equation}
  \hat{y}_i = \operatorname{arg\,max}_j \operatorname{softmax}(W\mathbf{x}_i + \mathbf{b})_j
\end{equation}

\textbf{Training configuration:} 3 epochs, batch size 8, learning rate $2 \times 10^{-5}$, max sequence length 128, AdamW optimiser, unweighted cross-entropy loss.

\subsection{Model 3: Multilingual E5-base (embedding model with token classifier)}
\label{sec:me5}

Multilingual-E5-base (\texttt{intfloat/multilingual-e5-base})~\cite{wang2023e5} is an XLM-RoBERTa-based sentence embedding model ($\approx$278M parameters) trained using contrastive learning on $\approx$1 billion multilingual text pairs. A token classification head is attached for NER:
\begin{equation}
  Z = XW + b
\end{equation}
where $X \in \mathbb{R}^{N \times 768}$ is the matrix of contextual token embeddings and $Z \in \mathbb{R}^{N \times 13}$ contains output logits.

\textbf{Class weights} are computed as:
\begin{equation}
  w_c = \frac{\sum_{j \in \text{non-O}} \text{count}_j}{\lvert\text{non-O}\rvert \cdot \text{count}_c}
\end{equation}
clipped to $[0.05,\; 8.0]$ for numerical stability.

\textbf{Training configuration:} 15 epochs, batch size 8 (train) / 16 (eval), gradient accumulation 2 steps, learning rate $2 \times 10^{-5}$, AdamW (weight decay 0.01, warmup ratio 0.1), seed 42.

\subsection{Model 4: Qwen2.5 (LLM via Ollama, few-shot prompting)}
\label{sec:llm}

Qwen2.5~\cite{yang2024qwen} is Alibaba's open-source 7B-parameter instruction-tuned LLM supporting 29+ languages including Russian and Kazakh. The model was deployed locally via the \textit{Ollama} serving framework, enabling inference without external API calls.

\textbf{Prompt construction.} The full inference prompt is constructed as:
\begin{equation}
  \text{prompt}_i = \text{SYS} \oplus E_{1..12} \oplus x_i
\end{equation}
where SYS is a system instruction defining the 6 entity types and extraction rules for Kazakh/Russian transit complaints, $E_{1..12}$ are 12 balanced in-context examples (2 per entity type) from the training set, and $x_i$ is the raw complaint text.

\textbf{Post-filtering.} To remove hallucinated spans, a post-filter is applied:
\begin{equation}
  \hat{\mathcal{E}}_i^{(c)} = \left\{v \in \hat{\mathcal{E}}_i^{(c)} \mid |v| \geq 2,\; v \notin \mathcal{W}_{\text{stop}},\; \phi(x_i, c)\right\}
\end{equation}
where $|v|$ is the character length, $\mathcal{W}_{\text{stop}}$ is a set of noise words, and $\phi(x_i, c)$ is a keyword check required for \texttt{STREET\_INTERSECTION} entities.

\textbf{Inference configuration:} max tokens 300, timeout 90\,s / 3 retries, seed 42, 12 few-shot examples. Evaluated on the first 300 test samples due to computational constraints.

\section{Baseline methods for comparison}
\label{sec:baselines}

Four baselines were implemented to contextualise the performance of the main models.

\subsection{Regex-based rule extraction}
\label{sec:regex_baseline}

A deterministic rule-based extractor using regular expressions was implemented for the three structurally regular entity types: \texttt{ROUTE\_NUM}, \texttt{BUS\_ID}, and \texttt{PLATE\_NUM}. Patterns were hand-crafted from a manual inspection of 100 training examples per entity type. Route numbers were matched by numeric sequences optionally preceded by Kazakh/Russian keywords (\foreignlanguage{russian}{маршрут}, \foreignlanguage{russian}{бағыт}, №). Plate numbers were matched by the standard Kazakhstani licence plate format. The regex baseline serves as a lower-bound reference for structured entity types.

\subsection{Frozen XLM-R linear probe}
\label{sec:frozen_baseline}

A linear probe was trained on frozen XLM-RoBERTa representations: the transformer backbone weights were fixed at their pre-trained values, and only the token classification head $W \in \mathbb{R}^{768 \times K}$ was trained. This baseline isolates the contribution of fine-tuning (weight updates to the backbone) from the contribution of the pre-trained contextual representations alone.

\subsection{CRF baseline}
\label{sec:crf_baseline}

A Conditional Random Field (CRF) model was trained on handcrafted token-level features: word form, prefix/suffix patterns (length 2--4), capitalisation flags, digit flags, and a domain gazetteer of known stop names and street names from the Astana city registry. CRF training used L-BFGS optimisation with L1 + L2 regularisation. The CRF baseline represents the statistical (pre-deep-learning) state of the art.

\subsection{Most-frequent-class baseline}
\label{sec:mfc_baseline}

A trivial baseline that assigns the $O$ (outside) label to every token, reflecting the dominant class in a heavily imbalanced corpus. This baseline defines the floor of macro-averaged F1: because all entity classes receive F1\,=\,0.000 under this scheme, the macro F1 is 0.000 regardless of the accuracy on $O$-class tokens.

\section{NER quality assessment metrics (precision, recall, F1-score, confusion matrix, boundary stability)}
\label{sec:metrics}

All models are evaluated using standard span-level NER metrics computed by the \texttt{seqeval} library, which requires exact boundary and label match for a true positive:
\begin{equation}
  \text{Precision} = \frac{|\text{TP}|}{|\text{TP}| + |\text{FP}|}, \quad
  \text{Recall} = \frac{|\text{TP}|}{|\text{TP}| + |\text{FN}|}, \quad
  F_1 = 2 \cdot \frac{\text{Precision} \cdot \text{Recall}}{\text{Precision} + \text{Recall}}
  \label{eq:f1_ch2}
\end{equation}

Macro-averaged F1 weights each entity type equally, regardless of frequency, making it sensitive to performance on rare classes.

\textbf{Confusion matrix.} A token-level confusion matrix across all 14 label classes (13 BIO entity labels + $O$) is computed to identify systematic label collapse patterns (e.g., a model that predominantly predicts \texttt{B-STOP\_NAME} for all entity tokens).

\textbf{Boundary stability.} For each predicted entity span, boundary stability measures the fraction of predictions where both the left and right boundaries are correct, given that the label is correct. A model with high precision but low boundary stability produces correctly labelled but incorrectly delimited spans, which degrades downstream geocoding accuracy.
